\documentclass[../../main.tex]{subfiles}

\begin{document}

\section{可数可加性、连续性以及 Borel-Cantelli 引理}\label{section:ch02-可数可加}

\begin{definition}[Lebesgue 测度]\label{definition:ch02-Lebesgue测度}
外测度这一集函数在可测集类上的限制称为 Lebesgue 测度. 它记为 $m$, 因此若 $E$ 是可测集, 它的 Lebesgue 测度 $m(E)$ 定义为
\begin{align*}
m(E)=m^{*}(E).
\end{align*}
\end{definition}

下面的命题十分重要.

\begin{proposition}\label{proposition:ch02-可数可加}
Lebesgue 测度是可数可加的, 即若 $\{E_{k}\}_{k=1}^{\infty}$ 是可测集的可数不交族, 则它的并集 $\bigcup_{k=1}^{\infty}E_{k}$ 也是可测的且
\begin{align*}
m\left(\bigcup_{k=1}^{\infty}E_{k}\right)=\sum_{k=1}^{\infty}m(E_{k}).
\end{align*}
\end{proposition}

\begin{proof}
\Cref{proposition:ch02-可数并可测} 告诉我们 $\bigcup_{k=1}^{\infty}E_{k}$ 是可测的. 根据\cref{proposition:ch02-外测度次可加}, 外测度是可数次可加的. 因此
\begin{align}
m\left(\bigcup_{k=1}^{\infty}E_{k}\right)\leqslant\sum_{k=1}^{\infty}m(E_{k}).\label{align:ch02-上和}
\end{align}
剩下来要证明反向的不等式. 根据\cref{proposition:ch02-有限不交可加}, 对每个自然数 $n$,
\begin{align*}
m\left(\bigcup_{k=1}^{n}E_{k}\right)=\sum_{k=1}^{n}m(E_{k}).
\end{align*}
由于 $\bigcup_{k=1}^{\infty}E_{k}$ 包含 $\bigcup_{k=1}^{n}E_{k}$, 根据外测度的单调性和前一个等式, 对每个 $n$,
\begin{align*}
m\left(\bigcup_{k=1}^{\infty}E_{k}\right)\geqslant\sum_{k=1}^{n}m(E_{k}).
\end{align*}
该不等式的左边与 $n$ 无关. 因此,
\begin{align}
m\left(\bigcup_{k=1}^{\infty}E_{k}\right)\geqslant\sum_{k=1}^{\infty}m(E_{k}).\label{align:ch02-下和}
\end{align}
从不等式\eqref{align:ch02-上和}和\eqref{align:ch02-下和}得出这些是等式.
\end{proof}

根据\cref{proposition:ch02-区间外测度}, 区间的外测度是它的长度, 且根据\cref{proposition:ch02-外测度平移}, 外测度是平移不变的. 因此前一个命题完成了以下定理的证明, 它是本章的主要目标.

\begin{theorem}\label{theorem:ch02-测度定理}
定义在 Lebesgue 可测集的 $\sigma$ 代数上的集函数 Lebesgue 测度, 赋予任何区间它的长度, 是平移不变且可数可加的.
\end{theorem}

集合的可数族 $\{E_{k}\}_{k=1}^{\infty}$ 称为是上升的, 若对每个 $k$, $E_{k}\subseteq E_{k+1}$; 而称为是下降的, 若对每个 $k$, $E_{k+1}\subseteq E_{k}$.

\begin{theorem}[测度的连续性]\label{theorem:ch02-测度连续}
Lebesgue 测度具有以下连续性质:
\begin{enumerate}[(i)]
\item 若 $\{A_{k}\}_{k=1}^{\infty}$ 是上升的可测集族, 则
\begin{align}
m\left(\bigcup_{k=1}^{\infty}A_{k}\right)=\lim_{k\to\infty}m(A_{k}).\label{align:ch02-上升连续}
\end{align}
\item 若 $\{B_{k}\}_{k=1}^{\infty}$ 是下降的可测集族且 $m(B_{1})<\infty$, 则
\begin{align}
m\left(\bigcap_{k=1}^{\infty}B_{k}\right)=\lim_{k\to\infty}m(B_{k}).\label{align:ch02-下降连续}
\end{align}
\end{enumerate}
\end{theorem}

\begin{proof}
我们首先证明 (i). 若存在指标 $k_{0}$ 使得 $m(A_{k_{0}})=\infty$, 则根据外测度的单调性, $m(\bigcup_{k=1}^{\infty}A_{k})=\infty$, 且对所有 $k\geqslant k_{0}$, $m(A_{k})=\infty$. 因此\eqref{align:ch02-上升连续}成立, 这是因为等式的两边都等于 $\infty$. 剩下来要考虑对所有 $k$, $m(A_{k})<\infty$ 的情形. 定义 $A_{0}=\varnothing$, 接着对每个 $k\geqslant 1$ 定义 $C_{k}=A_{k}\sim A_{k-1}$. 根据构造, 由于序列 $\{A_{k}\}_{k=1}^{\infty}$ 是上升的,
\begin{align*}
\{C_{k}\}_{k=1}^{\infty}\text{ 是不交的 且 }\bigcup_{k=1}^{\infty}A_{k}=\bigcup_{k=1}^{\infty}C_{k}.
\end{align*}
根据 $m$ 的可数可加性,
\begin{align}
m\left(\bigcup_{k=1}^{\infty}A_{k}\right)=m\left(\bigcup_{k=1}^{\infty}C_{k}\right)=\sum_{k=1}^{\infty}m(A_{k}\sim A_{k-1}).\label{align:ch02-连续分解}
\end{align}
由于 $\{A_{k}\}_{k=1}^{\infty}$ 是上升的, 我们从测度的分割性质推出
\begin{align*}
\sum_{k=1}^{\infty}m(A_{k}\sim A_{k-1})&=\sum_{k=1}^{\infty}\left[m(A_{k})-m(A_{k-1})\right]\\
&=\lim_{n\to\infty}\sum_{k=1}^{n}\left[m(A_{k})-m(A_{k-1})\right]\\
&=\lim_{n\to\infty}\left[m(A_{n})-m(A_{0})\right].
\end{align*}
由于 $m(A_{0})=m(\varnothing)=0$, \eqref{align:ch02-上升连续}从\eqref{align:ch02-连续分解}得出.

为证明 (ii), 对每个 $k$ 我们定义 $D_{k}=B_{1}\sim B_{k}$. 由于序列 $\{B_{k}\}_{k=1}^{\infty}$ 是下降的, 序列 $\{D_{k}\}_{k=1}^{\infty}$ 是上升的. 根据 (i) 部分,
\begin{align*}
m\left(\bigcup_{k=1}^{\infty}D_{k}\right)=\lim_{k\to\infty}m(D_{k}).
\end{align*}
根据 De Morgan 恒等式,
\begin{align*}
\bigcup_{k=1}^{\infty}D_{k}=\bigcup_{k=1}^{\infty}\left[B_{1}\sim B_{k}\right]=B_{1}\sim\bigcap_{k=1}^{\infty}B_{k}.
\end{align*}
另一方面, 根据测度的分割性质, 对每个 $k$, 由于 $m(B_{k})<\infty$, $m(D_{k})=m(B_{1})-m(B_{k})$. 因此
\begin{align*}
m\left(B_{1}\sim\bigcap_{k=1}^{\infty}B_{k}\right)=\lim_{n\to\infty}\left[m(B_{1})-m(B_{n})\right].
\end{align*}
再一次用分割性质我们得到等式\eqref{align:ch02-下降连续}.
\end{proof}

对于可测集 $E$, 我们说一个性质在 $E$ 上几乎处处成立, 或它对几乎所有 $x\in E$ 成立, 若存在 $E$ 的子集 $E_{0}$ 使得 $m(E_{0})=0$, 且该性质对所有 $x\in E\sim E_{0}$ 成立.

\begin{theorem}[Borel-Cantelli 引理]\label{theorem:ch02-BorelCantelli}
令 $\{E_{k}\}_{k=1}^{\infty}$ 为满足 $\sum_{k=1}^{\infty}m(E_{k})<\infty$ 的可测集的可数族. 则几乎所有 $x\in\mathbf{R}$ 属于至多有限多个 $E_{k}$.
\end{theorem}

\begin{proof}
对每个 $n$, 根据 $m$ 的可数次可加性,
\begin{align*}
m\left(\bigcup_{k=n}^{\infty}E_{k}\right)\leqslant\sum_{k=n}^{\infty}m(E_{k})<\infty.
\end{align*}
因此, 根据测度的连续性,
\begin{align*}
m\left(\bigcap_{n=1}^{\infty}\left[\bigcup_{k=n}^{\infty}E_{k}\right]\right)=\lim_{n\to\infty}m\left(\bigcup_{k=n}^{\infty}E_{k}\right)\leqslant\lim_{n\to\infty}\sum_{k=n}^{\infty}m(E_{k})=0.
\end{align*}
因此几乎所有 $x\in\mathbf{R}$ 不属于 $\bigcap_{n=1}^{\infty}\left[\bigcup_{k=n}^{\infty}E_{k}\right]$, 因而属于至多有限多个 $E_{k}$.
\end{proof}

集函数 Lebesgue 测度承袭了 Lebesgue 外测度具有的性质. 为了将来的参考我们命名这些性质中的一些性质.

\textbf{(有限可加性)} 对任何可测集的有限不交族 $\{E_{k}\}_{k=1}^{n}$,
\begin{align*}
m\left(\bigcup_{k=1}^{n}E_{k}\right)=\sum_{k=1}^{n}m(E_{k}).
\end{align*}

\textbf{(单调性)} 若 $A$ 和 $B$ 是可测集且 $A\subseteq B$, 则
\begin{align*}
m(A)\leqslant m(B).
\end{align*}

\textbf{(分割性)} 若进一步, $A\subseteq B$ 且 $m(A)<\infty$, 则
\begin{align*}
m(B\sim A)=m(B)-m(A).
\end{align*}
因此若 $m(A)=0$, 则
\begin{align*}
m(B\sim A)=m(B).
\end{align*}

\textbf{(可数单调性)} 对任何覆盖可测集 $E$ 的可测集的可数族 $\{E_{k}\}_{k=1}^{\infty}$,
\begin{align*}
m(E)\leqslant\sum_{k=1}^{\infty}m(E_{k}).
\end{align*}
可数单调性是测度的单调性和可数次可加性的合并. 它经常被使用.

\begin{remark}
在将要进行的 Lebesgue 积分的研究中, 事实是显然的: 相比于 Riemann 积分, Lebesgue 测度的可数可加性使得 Lebesgue 积分具有决定性的优势.
\end{remark}

\begin{exercise}\label{exercise:ch02-24}
证明: 若 $E_{1}$ 和 $E_{2}$ 是可测的, 则
\begin{align*}
m(E_{1}\cup E_{2})+m(E_{1}\cap E_{2})=m(E_{1})+m(E_{2}).
\end{align*}
\end{exercise}

\begin{exercise}\label{exercise:ch02-25}
证明在关于测度的连续性的定理的 (ii) 部分假设 $m(B_{1})<\infty$ 是必要的.
\end{exercise}

\begin{exercise}\label{exercise:ch02-26}
令 $\{E_{k}\}_{k=1}^{\infty}$ 为可测集的可数不交族. 证明对任意集合 $A$,
\begin{align*}
m^{*}\left(A\cap\bigcup_{k=1}^{\infty}E_{k}\right)=\sum_{k=1}^{\infty}m^{*}(A\cap E_{k}).
\end{align*}
\end{exercise}

\begin{exercise}\label{exercise:ch02-27}
令 $\mathcal{M}'$ 为 $\mathbf{R}$ 的子集的 $\sigma$ 代数, 而 $m'$ 是 $\mathcal{M}$ 上的集函数, 它在 $[0,\infty]$ 取值, 是可数可加的, 且使得 $m'(\varnothing)=0$.
\begin{enumerate}[(i)]
\item 证明 $m'$ 是有限可加的、单调的、可数单调的且具有分割性质.
\item 证明 $m'$ 具有与 Lebesgue 测度相同的连续性质.
\end{enumerate}
\end{exercise}

\begin{exercise}\label{exercise:ch02-28}
证明测度的连续性与测度的有限可加性蕴涵测度的可数可加性.
\end{exercise}

\end{document}